\section{Discussion}
\label{sec:discussion}

\subsection{Mechanisms}
\label{sec:disc-mechanism}
Reading fabric through the three registers of type, measurement and mechanism does
analytical work here. The provisional planning types are recovered as measurable
form (Table~\ref{tab:fabric-profiles}) but consolidate into four robust
morphometric super-types, and the mechanism layer then explains why the
hottest of these is the coarse-grain industrial fabric: large, contiguous
impervious footprints, the leading morphological term in the SHAP attribution. At
the same time, at metropolitan scale urban form is a secondary control on summer
land-surface temperature once the coastal gradient is taken into account.
Morphology alone carries modest but genuine skill
(spatial-block $R^2 \approx 0.09$), roughly tripling to $R^2 \approx 0.27$ once
coastal and topographic context is admitted. This is consistent with marginal/interaction
accounts that find form effects real but conditional \citep{wangZ2025MarginalInteraction},
and it cautions against importing the strong morphology-temperature coefficients
reported for inland or fine-scale settings \citep{liu2026Nonlinear, li2026BlockGrid}
into a maritime metropolis without the coastal control.

\label{sec:disc-positioning}
The contribution is the binding, in one commensurable unit, of reproducible
morphometrics, measured heat, accessibility and social vulnerability, and their
conversion into trade-off-aware priorities. Relative to descriptive typologies
\citep{araldi2025MultiLevel, fleischmann2022MethodolFoundati} the workflow attaches
hazard and equity outcomes to fabric; relative to grid/LCZ-scale SHAP-LST studies
\citep{wangZ2025PlainPlateau, li2026BlockGrid} it operates at the street/cell scale
and feeds attribution into a decision step; and relative to coarse vulnerability
indices \citep{turner2025Guhvi} it ties vulnerability to the form that produces
differential exposure. The two stages are complementary rather than nested: the SHAP attributions diagnose
why a fabric is hot and so guide the type of intervention each priority
fabric needs, while the optimizer ranks where to act using measured
land-surface temperature rather than a SHAP-derived heat axis. The sharpest
illustration is an inversion the integrated unit makes visible: the waterfront
fabric is the coolest yet the highest-priority, because it is the most
demographically vulnerable (elderly share 19.0\%) and the most coastally exposed.
The coastal-exposure driver is partly definitional, so the claim is best read as:
a cool fabric ranks first once vulnerability and exposure are
admitted. A heat-only assessment would have mis-ranked it, providing a concrete instance of
the blind spot the introduction identifies. Although the main-text vulnerability
axis is an age-structure proxy, robustness tests re-run the priority pipeline with
two age-independent axes derived from open GHSL data: a building-age proxy
(thermal envelope and pre-code construction) and a dwelling-crowding proxy. This
keeps the waterfront top-ranked under building age and under a three-axis
composite, and in the top tier under crowding (Appendix~\ref{app:vuln-robust},
Table~\ref{tab:vuln-robust}). The building-age ranking reproduces the age-based ordering
almost exactly (Spearman $\rho = 0.96$), so the paradox is a property of the coupled
need structure rather than of the demographic operationalization of vulnerability.

\subsection{Policy translation}
\label{sec:disc-policy}
For \.{I}zmir, the priorities translate fairly directly. The high-high need clusters
(historic core, waterfront, industrial) are where shading, greening and
permeable-surface interventions would couple heat and equity gains
\citep{neumann2026ClimateJustice}; these locations represent critical targets for
ecosystem-based adaptation and targeted nature-based solutions to build urban
resilience \citep{geneletti2016EcosysteBased, wamsler2020EnvironmClimate,
babi2021NexusBetween}. The industrial fabric's footprint-driven heat
points to roof/yard greening and albedo measures, drawing on documented tree and green
infrastructure cooling benefits that optimize population-exposure reduction
\citep{yu2020CriticalReview, schwaab2021RoleUrban, marando2022UrbanHeat,
massaro2023SpatiallOptimize}. The waterfront's priority is driven by vulnerability
and exposure rather than heat, implying social and evacuation-access measures over
cooling alone. Because the spatial-statistics layer localizes need (LISA, $G_i^{*}$)
and the priorities are reported with their weight-robustness, the output is usable
as a transparent screening input to municipal climate-action planning rather than
a black-box score.

Table~\ref{tab:policy-translation} converts the mechanism and ranking outputs
into intervention packages. The SHAP values are diagnostic associations rather
than intervention effects: for example, peripheral expansion carries a mean
green-cover attribution of $-0.81\,\si{\celsius}$, whereas the large-footprint
term in industrial/logistics carries $+0.40\,\si{\celsius}$. This
$1.21\,\si{\celsius}$ model-attributable contrast indicates where additional
green cover and large-roof retrofits align with the fitted heat mechanism; it is
not a claim that either intervention would cause the full contrast.

\begin{table}[htbp]
  \centering\footnotesize
  \caption{Policy translation by fabric stratum. Dominant mechanisms are the
  largest absolute stratum-mean SHAP terms in Figure~\ref{fig:shap}; priority is
  the Monte-Carlo mean rank from Table~\ref{tab:priority} (1 = highest).}
  \label{tab:policy-translation}
  \begin{tabular}{@{}p{0.18\textwidth}p{0.24\textwidth}p{0.10\textwidth}p{0.38\textwidth}@{}}
    \toprule
    Fabric stratum & Dominant SHAP mechanism & Priority rank & Recommended intervention package \\
    \midrule
    Waterfront transformation & Coastal proximity cooling ($-1.84\,\si{\celsius}$) & 1.6 & Evacuation access, heat-health outreach, shaded assembly areas and coastal-exposure planning. \\
    Historic core & Green-cover deficit ($+0.32\,\si{\celsius}$) & 2.7 & Courtyard and pocket greening, reversible shade, cool roofs and protected pedestrian routes. \\
    Peripheral expansion & Green cover cooling ($-0.81\,\si{\celsius}$) & 3.0 & Preserve existing green cover, add street trees and connect daily services and assembly access. \\
    Industrial / logistics & Large footprints warming ($+0.40\,\si{\celsius}$) & 4.2 & Cool and green roofs, yard trees, high-albedo paving and worker shade. \\
    Apartment block & Green-cover deficit ($+0.32\,\si{\celsius}$) & 5.1 & Courtyard trees, facade and roof greening, and shaded walking links. \\
    Hillside / incremental & Slope context ($-0.29\,\si{\celsius}$) & 5.7 & Slope-safe evacuation routes, distributed shade and small-scale permeable surfaces. \\
    Grid residential & Green cover cooling ($-0.20\,\si{\celsius}$) & 5.8 & Retain canopy, shade intersections and improve access to services and assembly areas. \\
    \bottomrule
  \end{tabular}
\end{table}

\subsection{Limitations}
\label{sec:disc-limits}
Several caveats bound these claims. (i) The fabric strata are \emph{provisional} and
rule-based; they are earmarked for refinement with planning documents and
urban-design expertise, and the clustering already signals that seven types
over-partition a four-mode morphometric space. (ii) The heat layer is screening,
not prediction: a modest, conservatively reported $R^2$ and a coast-dominated signal mean
the SHAP mechanism attributions describe associations within this sample, not a
transferable temperature model. The attributions are nonetheless spatially stable: the
feature-importance ordering holds across spatial-block folds (Spearman $= 0.86$) and
every high-importance predictor keeps its sign of effect in all folds
(Appendix~\ref{app:shap-stability}), so the mechanism reading does not hinge on any
single region. (iii) The analysis uses a multi-year summer-mean
LST over the full census of 3{,}777 cells; seasonal/diurnal contrasts and a fine-scale
(pixel) attribution are deferred to future work. (iv) The main-text social-vulnerability axis rests on
age-structure proxies (elderly share and age dependency from the 2024 ADNKS register);
over the full census every grid cell falls within a populated \emph{mahalle}, so no
imputation or spatial fallback is required. We show the equity inversion is robust to
this operationalization by re-deriving vulnerability from two age-independent GHSL axes
(building age and dwelling crowding) and confirming the waterfront stays top-tier
(Appendix~\ref{app:vuln-robust}); nonetheless, richer socioeconomic inputs (income,
tenure, health-service access) would further sharpen the equity axis. (v) The five-axis Pareto
screen is uninformative at the fabric level (no fabric dominated); applied at the
cell level it is selective (223 of 3{,}777 cells non-dominated), which is the level at
which the dominance result should be read. (vi) Built-form intensity uses
footprint area and measured floor counts; hazard treatment is non-hydrodynamic.
(vii) Pluvial susceptibility is provided only as a topographic screening overlay
and does not enter the five priority axes. Hydrodynamic validation is required
before it can be promoted to a ranked need axis. (viii) The tessellation-as-plot-proxy assumption was not validated against cadastral
or plan records for \.{I}zmir in this study; such a check is left for the
refinement stage.

\subsection{Outlook}
\label{sec:disc-transfer}
The workflow runs on open and openly licensed inputs, including OpenStreetMap, ESA
WorldCover, Landsat via Google Earth Engine, an open DEM and the public population
register. Official municipal building and street layers supplement these sources and
materially improved completeness over crowd-sourced footprints. With frozen
software versions, logged parameters and released code, the same chain transfers to
other coastal metropolitan regions, where the coastal-gradient finding suggests
that context terms must travel with the morphological ones. Concrete next
candidates include Mediterranean coastal cities with comparable open-data coverage,
such as Barcelona, Marseille, Athens, Antalya and Alexandria, all of which combine
dense historical cores, post-war expansion rings and littoral exposure gradients
analogous to \.{I}zmir's. In each case the core inputs include OpenStreetMap building
footprints, Copernicus Sentinel-2/Landsat surface temperature, ESA WorldCover land
use and national census or population-register data. These are publicly available,
making the transfer operationally feasible. The critical lesson from the \.{I}zmir
census is that, in maritime cities, contextual drivers (the coastal gradient) must
travel with the morphological ones; any replication that omits the context control
risks over-attributing thermal variance to built form.

\label{sec:disc-future}
Several extensions are planned for the full study. First, the current multi-year
summer-mean LST will be decomposed into seasonal and diurnal contrasts (day versus
night, summer versus winter) to reveal fabric-specific thermal behaviour that a
single composite obscures \citep{wangZ2025MarginalInteraction}. Second, the SHAP
attribution will be pushed from the cell level to the pixel scale, sharpening the
morphology-heat link and enabling within-cell identification of the built-form
elements that most amplify or mitigate surface temperature. Third, multiscale
geographically weighted regression (MGWR) will be introduced to capture local
non-stationarity in the heat-morphology relationship, extending the multiscale
Istanbul evidence base with spatially varying coefficients
\citep{zhao2018GeographWeighted, eminoglu2026Diagnosing}. Fourth, the
screening-level pluvial-exposure
layer will be replaced by, or validated against, hydrodynamic flood modelling to
move from topographic susceptibility to physically based inundation depths
\citep{imroz2025IntegratedAssess}. Finally, deployment of the workflow on a second
coastal city will provide an inter-city transferability test and determine whether
the coastal-context finding generalizes beyond the Aegean setting.
